\documentclass[12pt, a4paper]{article}

% 1. Cấu hình Tiếng Việt và Font Times New Roman 13pt chuẩn
\usepackage[utf8]{inputenc}
\usepackage[T5]{fontenc}
\usepackage[vietnamese]{babel}
\usepackage{newtxtext,newtxmath}

\makeatletter
\renewcommand{\normalsize}{\@setfontsize\normalsize{13pt}{16pt}}
\makeatother
\normalsize

% 2. Gói Toán học & Ký hiệu bổ trợ
\usepackage{amsmath, amssymb, amsfonts, mathrsfs, amsthm}
\usepackage{graphicx}
\usepackage{fontawesome}
\usepackage{booktabs, tabularx, array, multirow}
\usepackage{enumitem}

% 3. Đồ họa TikZ, Bảng biến thiên & Biểu đồ
\usepackage{tikz}
\usetikzlibrary{calc, intersections, angles, quotes, patterns, arrows.meta, 3d}
\usepackage{pgfplots}
\pgfplotsset{compat=1.18}
\usepackage{tkz-euclide}
\usepackage{tkz-tab}

\renewcommand{\vec}[1]{\overrightarrow{#1}}

% 4. Hộp sư phạm (Tcolorbox)
\usepackage[many]{tcolorbox}
\tcbuselibrary{skins, breakable}

\newtcolorbox{pedabox}[1][]{
    enhanced,
    colback=blue!5!white,
    colframe=blue!75!black,
    fonttitle=\bfseries\sffamily,
    title=#1,
    boxrule=0.8pt,
    arc=2mm,
    breakable,
    left=3mm, right=3mm, top=2mm, bottom=2mm
}

\newtcolorbox{scaffoldbox}[1][]{
    enhanced,
    colback=orange!5!white,
    colframe=orange!85!black,
    fonttitle=\bfseries\sffamily,
    title=#1,
    boxrule=0.8pt,
    arc=2mm,
    breakable,
    left=3mm, right=3mm, top=2mm, bottom=2mm
}

\newtcolorbox{aibox}[1][]{
    enhanced,
    colback=green!5!white,
    colframe=teal!80!black,
    fonttitle=\bfseries\sffamily,
    title=#1,
    boxrule=0.8pt,
    arc=2mm,
    breakable,
    left=3mm, right=3mm, top=2mm, bottom=2mm
}

% 5. Căn lề chuẩn văn bản giáo án
\usepackage{geometry}
\geometry{
    a4paper,
    left=25mm,
    right=15mm,
    top=20mm,
    bottom=20mm
}

% 6. Header / Footer chuẩn hóa (BẮT BUỘC CHÍNH XÁC NỘI DUNG NÀY)
\usepackage{fancyhdr}
\usepackage{lastpage}
\pagestyle{fancy}
\fancyhf{}

\fancyhead[L]{\small \textit{Kế hoạch bài dạy Toán 11}}
\fancyhead[R]{\textbf{Trang \thepage/\pageref{LastPage}}}
\fancyfoot[L]{\small \textit{Tổ Toán - Tin}}
\fancyfoot[R]{\small \textit{Trường THCS\&THPT Hồng Vân}}
\renewcommand{\headrulewidth}{0.4pt}
\renewcommand{\footrulewidth}{0.4pt}

% 7. Giãn dòng & Giãn đoạn theo chuẩn Nghị định 30/2020/NĐ-CP
\usepackage{setspace}
\setstretch{1.15}
\setlength{\parindent}{1.0cm}
\setlength{\parskip}{5pt}

\begin{document}

\begin{center}
    {\Large \textbf{KẾ HOẠCH BÀI DẠY MÔN TOÁN LỚP 11}}\\[0.4em]
    {\LARGE \textbf{BÀI 8. MẪU SỐ LIỆU GHÉP NHÓM}}\\[0.5em]
    \textbf{Môn học:} Toán 11 | \textbf{Thời lượng:} 01 tiết (Tiết 18 theo PPCT)
\end{center}

\vspace{0.5em}
\hrule
\vspace{1em}

\section*{I. MỤC TIÊU DẠY HỌC}

\subsection*{1. Kiến thức, Năng lực}

\textbf{a) Năng lực Toán học đặc thù (nguyên văn chuẩn CT GDPT 2018):}
\begin{itemize}[leftmargin=1.5em]
    \item \textbf{Năng lực tư duy và lập luận toán học:}
    \begin{itemize}
        \item Nhận biết và giải thích được lí do cần thiết phải ghép nhóm các số liệu rời rạc thành các khoảng/nửa khoảng khi xử lý các tập dữ liệu có quy mô lớn hoặc mang tính chất biến số liên tục (như chiều cao, cân nặng, thời gian, doanh thu).
        \item Hiểu rõ và phát biểu chính xác khái niệm nhóm số liệu $[a_i; a_{i+1})$, đầu mút trái $a_i$, đầu mút phải $a_{i+1}$, độ dài nhóm $h = a_{i+1} - a_i$ và giá trị đại diện của nhóm $c_i = \dfrac{a_i + a_{i+1}}{2}$.
        \item Phân tích và giải thích được mối quan hệ giữa tần số tuyệt đối $m_i$ và tần số tương đối $f_i = \dfrac{m_i}{n} \cdot 100\%$ của từng nhóm.
    \end{itemize}
    \item \textbf{Năng lực giải quyết vấn đề toán học:}
    \begin{itemize}
        \item Thiết lập được bảng tần số ghép nhóm và bảng tần số tương đối ghép nhóm từ một mẫu số liệu ban đầu chưa ghép nhóm.
        \item Xác định chính xác các giá trị quan sát thuộc vào từng nhóm theo quy tắc nửa khoảng $[a_i; a_{i+1})$ (giá trị $x$ thỏa mãn $a_i \le x < a_{i+1}$, riêng nhóm cuối cùng thường lấy cả hai đầu mút $[a_k; a_{k+1}]$).
        \item Thiết lập và vẽ được biểu đồ tần số tương đối ghép nhóm (Histogram) ở dạng biểu đồ cột liền kề hoặc biểu đồ đoạn thẳng.
    \end{itemize}
    \item \textbf{Năng lực mô hình hóa toán học:}
    \begin{itemize}
        \item Vận dụng kĩ thuật ghép nhóm số liệu để khảo sát, mô hình hóa các hiện tượng thực tế trong trường học và đời sống: khảo sát chiều cao học sinh khối 11, thống kê thời gian ôn bài tại nhà, phân tích mức tiêu thụ điện năng của các hộ gia đình.
    \end{itemize}
    \item \textbf{Năng lực giao tiếp toán học:}
    \begin{itemize}
        \item Sử dụng chính xác các thuật ngữ thống kê: mẫu số liệu ghép nhóm, khoảng nửa mở, tần số, tần số tương đối, giá trị đại diện, biểu đồ Histogram; đọc hiểu và truyền đạt thông tin từ bảng số liệu ghép nhóm đến bạn học.
    \end{itemize}
    \item \textbf{Năng lực sử dụng công cụ, phương tiện học toán:}
    \begin{itemize}
        \item Khai thác máy tính cầm tay (MTCT Casio fx-580VN X) để tính nhanh tỉ lệ phần trăm và giá trị đại diện; sử dụng bảng tính điện tử (Microsoft Excel / Google Sheets) để xử lý dữ liệu và vẽ biểu đồ.
    \end{itemize}
\end{itemize}

\textbf{b) Năng lực số (theo Thông tư 02/2025/TT-BGDĐT):}
\begin{itemize}[leftmargin=1.5em]
    \item \textbf{Mã 5.1NC1b (Xử lý và phân tích dữ liệu bằng công cụ số):} Học sinh sử dụng phần mềm bảng tính Excel trên phòng máy 35 máy tính, áp dụng thành thạo hàm đếm mảng \texttt{FREQUENCY} để tự động phân nhóm và đếm tần số cho một tập dữ liệu gồm $50$ quan sát thực tế; tự động vẽ biểu đồ Histogram trực quan hóa sự phân bố của dữ liệu.
\end{itemize}

\textbf{c) Năng lực AI (theo Quyết định 2422/QĐ-BGDĐT):}
\begin{itemize}[leftmargin=1.5em]
    \item \textbf{Mã 11.C5.2 (Hiểu biết về thuật toán học máy và trí tuệ nhân tạo):} Học sinh nhận biết và nêu được sự tương đồng giữa việc ghép nhóm số liệu trong thống kê với thuật toán gom cụm (Clustering / K-Means) trong học máy (Machine Learning) của AI; hiểu rằng AI phân loại khách hàng, gom nhóm hành vi người dùng hay phân đoạn ảnh đều dựa trên nguyên lí phân chia khoảng cách dữ liệu tương tự như việc chia nhóm mẫu số liệu.
\end{itemize}

\textbf{d) Năng lực chung \& Phẩm chất:}
\begin{itemize}[leftmargin=1.5em]
    \item \textbf{Năng lực chung:} Tự chủ và tự học (tự thu thập số liệu thực tế cá nhân); giao tiếp và hợp tác (chia nhóm thảo luận phân chia khoảng cách nhóm hợp lí).
    \item \textbf{Phẩm chất:} Trung thực, khách quan khi thu thập và xử lý số liệu; cẩn thận, tỉ mỉ tránh đếm sót hoặc đếm lặp các giá trị biên.
\end{itemize}

\section*{II. THIẾT BỊ DẠY HỌC VÀ HỌC LIỆU}
\begin{itemize}[leftmargin=1.5em]
    \item \textbf{Giáo viên:} Kế hoạch bài dạy chi tiết chuẩn CV 5512; bài giảng PowerPoint; tệp bảng tính mẫu Excel chứa số liệu khảo sát thực tế và hàm \texttt{FREQUENCY}; phòng máy vi tính 35 máy; Phiếu học tập số 1 và số 2.
    \item \textbf{Học sinh:} Vở ghi bài, SGK Toán 11; thước kẻ, bút màu; máy tính cầm tay Casio fx-580VN X.
\end{itemize}

\section*{III. TIẾN TRÌNH DẠY HỌC (01 TIẾT --- 45 PHÚT)}

\begin{center}
    \framebox[1.0\textwidth]{\parbox{0.96\textwidth}{
        \textbf{CẤU TRÚC TIẾN TRÌNH TIẾT 18 (BÀI 8. MẪU SỐ LIỆU GHÉP NHÓM):}\\[0.3em]
        $\bullet$ \textbf{Hoạt động 1 (05 phút):} Khởi động --- Tình huống thực tế "Khảo sát thời gian tự học của học sinh lớp 11".\\[0.3em]
        $\bullet$ \textbf{Hoạt động 2 (22 phút):} Hình thành kiến thức mới:\\[0.2em]
        \quad - Mục 1: Khái niệm mẫu số liệu ghép nhóm và các yếu tố của nhóm (07 phút).\\[0.2em]
        \quad - Mục 2: Bảng tần số ghép nhóm \& Bảng tần số tương đối ghép nhóm (08 phút).\\[0.2em]
        \quad - Mục 3: Biểu đồ tần số tương đối ghép nhóm (Histogram) (04 phút).\\[0.2em]
        \quad - Mục 4: Tích hợp Năng lực số (Mã 5.1NC1b) \& Năng lực AI (Mã 11.C5.2) (03 phút).\\[0.3em]
        $\bullet$ \textbf{Hoạt động 3 (12 phút):} Luyện tập thiết lập bảng và vẽ biểu đồ.\\[0.3em]
        $\bullet$ \textbf{Hoạt động 4 (06 phút):} Vận dụng thực tiễn --- Phân tích dữ liệu điểm số thi thử tốt nghiệp.
    }}
\end{center}

\vspace{0.8em}
\hrule
\vspace{0.8em}

\subsubsection*{1. Hoạt động 1: Mở đầu / Khởi động (05 phút)}
\textbf{Chủ đề:} \textit{Làm thế nào để tổng hợp một danh sách số liệu đo lường dài dằng dặc?}
\begin{itemize}
    \item \textbf{a) Mục tiêu:} Tạo nhu cầu nhận thức về việc thu gọn và cấu trúc lại các tập dữ liệu định lượng lớn bằng phương pháp ghép nhóm thay vì giữ nguyên dãy số liệu rời rạc.
    \item \textbf{b) Nội dung:} Giáo viên chiếu bảng số liệu thô về thời gian sử dụng Internet trong một ngày (tính bằng phút) của $40$ học sinh lớp 11:
    \begin{center}
    \begin{tabular}{|cccccccc|}
    \hline
    $45$ & $60$ & $30$ & $75$ & $90$ & $120$ & $50$ & $65$ \\
    $80$ & $105$ & $40$ & $85$ & $70$ & $60$ & $95$ & $110$ \\
    $55$ & $60$ & $75$ & $130$ & $80$ & $90$ & $65$ & $85$ \\
    $100$ & $45$ & $70$ & $60$ & $85$ & $115$ & $90$ & $50$ \\
    $75$ & $80$ & $65$ & $95$ & $105$ & $60$ & $80$ & $70$ \\
    \hline
    \end{tabular}
    \end{center}
    GV đặt câu hỏi: Nhìn vào bảng số liệu trên, em có thể ngay lập tức trả lời được có bao nhiêu học sinh sử dụng Internet từ $60$ đến dưới $90$ phút không? Làm thế nào để người đọc dễ dàng quan sát phân bố của dữ liệu nhất?
    \item \textbf{c) Sản phẩm:} Học sinh nhận thấy việc đếm từng giá trị rất rối mắt và tốn thời gian; đề xuất ý tưởng gom các bạn có thời gian gần nhau vào từng nhóm (ví dụ nhóm dưới $60$ phút, nhóm $60 - 90$ phút, nhóm trên $90$ phút).
    \item \textbf{d) Tổ chức thực hiện:} GV chiếu bảng $\to$ HS quan sát, trao đổi nhanh $\to$ GV dẫn dắt: Trong thực tiễn thống kê, người ta ghép các số liệu vào các nhóm liên tiếp, gọi là \textbf{Mẫu số liệu ghép nhóm}.
\end{itemize}

\subsubsection*{2. Hoạt động 2: Hình thành kiến thức mới (22 phút)}

\paragraph{Mục 1: Mẫu số liệu ghép nhóm và các yếu tố của nhóm (07 phút)}
\begin{itemize}
    \item \textbf{Khái niệm:} Mẫu số liệu ghép nhóm là mẫu số liệu được cho dưới dạng bảng, trong đó các giá trị của mẫu được chia thành các nhóm liên tiếp, mỗi nhóm thường là một nửa khoảng có dạng $[a_i; a_{i+1})$ với $a_i < a_{i+1}$.
    \item \textbf{Các yếu tố cơ bản của một nhóm $[a_i; a_{i+1})$:}
    \begin{itemize}
        \item $a_i$ gọi là \textbf{đầu mút trái}, $a_{i+1}$ gọi là \textbf{đầu mút phải} của nhóm.
        \item $h = a_{i+1} - a_i$ gọi là \textbf{độ dài} của nhóm (thường chọn các nhóm có độ dài bằng nhau để dễ so sánh).
        \item $c_i = \dfrac{a_i + a_{i+1}}{2}$ gọi là \textbf{giá trị đại diện} của nhóm.
    \end{itemize}
    \item \textbf{Quy tắc xếp số liệu vào nhóm:} Giá trị $x$ thuộc vào nhóm $[a_i; a_{i+1})$ khi và chỉ khi:
    $$a_i \le x < a_{i+1}$$
    \textit{Lưu ý đặc biệt:} Riêng đối với nhóm cuối cùng $[a_k; a_{k+1}]$, ta thường lấy cả hai đầu mút (tức là $a_k \le x \le a_{k+1}$) để không bỏ sót giá trị lớn nhất của mẫu.
\end{itemize}

\begin{scaffoldbox}[Giàn giáo sư phạm cho HS TB - Yếu: Quy tắc "Trái lấy, Phải bỏ"]
    \textbf{Quy tắc nửa khoảng $[a; b)$ rất dễ nhầm lẫn:}\\
    $\bullet$ Dấu ngoặc vuông $[$ ở đầu mút trái $a$: \textbf{CÓ LẤY DẤU BẰNG} ($x \ge a$).\\
    $\bullet$ Dấu ngoặc tròn $)$ ở đầu mút phải $b$: \textbf{KHÔNG LẤY DẤU BẰNG} ($x < b$).\\
    $\bullet$ \textbf{Ví dụ kinh điển:} Cho nhóm $[60; 80)$.
    \begin{itemize}
        \item Số $60$: \textbf{THUỘC} nhóm $[60; 80)$.
        \item Số $79{,}9$: \textbf{THUỘC} nhóm $[60; 80)$.
        \item Số $80$: \textbf{KHÔNG THUỘC} nhóm $[60; 80)$ mà sẽ nhảy sang nhóm kế tiếp $[80; 100)$.
    \end{itemize}
    $\implies$ Giúp học sinh trung bình - yếu tuyệt đối không đếm trùng lặp hoặc bỏ sót số liệu!
\end{scaffoldbox}

\paragraph{Mục 2: Bảng tần số ghép nhóm \& Bảng tần số tương đối ghép nhóm (08 phút)}
\begin{itemize}
    \item \textbf{Tần số của nhóm ($m_i$):} Là số lượng các giá trị quan sát trong mẫu rơi vào nhóm $[a_i; a_{i+1})$.\\
    Tổng tần số của tất cả các nhóm bằng cỡ mẫu $n$:
    $$m_1 + m_2 + \dots + m_k = \sum_{i=1}^k m_i = n$$
    \item \textbf{Tần số tương đối của nhóm ($f_i$):} Là tỉ số phần trăm giữa tần số $m_i$ của nhóm đó và cỡ mẫu $n$:
    $$f_i = \dfrac{m_i}{n} \cdot 100\%$$
    Tổng tần số tương đối của tất cả các nhóm luôn bằng $100\%$:
    $$f_1 + f_2 + \dots + f_k = \sum_{i=1}^k f_i = 100\%$$
    \item \textbf{Cấu trúc tổng quát của Bảng tần số và tần số tương đối ghép nhóm:}
\end{itemize}

\begin{center}
\begin{tabularx}{0.95\linewidth}{|c|c|c|c|c|}
\hline
\textbf{Nhóm số liệu} & \textbf{Giá trị đại diện ($c_i$)} & \textbf{Tần số ($m_i$)} & \textbf{Tần số tương đối ($f_i$)} \\
\hline
$[a_1; a_2)$ & $c_1 = (a_1 + a_2)/2$ & $m_1$ & $f_1 = (m_1/n) \cdot 100\%$ \\
$[a_2; a_3)$ & $c_2 = (a_2 + a_3)/2$ & $m_2$ & $f_2 = (m_2/n) \cdot 100\%$ \\
\dots & \dots & \dots & \dots \\
$[a_k; a_{k+1}]$ & $c_k = (a_k + a_{k+1})/2$ & $m_k$ & $f_k = (m_k/n) \cdot 100\%$ \\
\hline
\textbf{Tổng} & --- & $n$ & $100\%$ \\
\hline
\end{tabularx}
\end{center}

\paragraph{Mục 3: Biểu đồ tần số tương đối ghép nhóm (Histogram) (04 phút)}
\begin{itemize}
    \item \textbf{Biểu đồ Histogram (biểu đồ cột liền kề):}
    \begin{itemize}
        \item Trục hoành nằm ngang biểu diễn các nhóm số liệu (ghi rõ các mốc đầu mút $a_1, a_2, \dots, a_{k+1}$).
        \item Trục tung thẳng đứng biểu diễn tần số tương đối $f_i$ (đơn vị $\%$).
        \item Các cột hình chữ nhật được dựng \textbf{sát liền kề nhau}, đáy mỗi cột là đoạn $[a_i; a_{i+1}]$ và chiều cao bằng $f_i$.
    \end{itemize}
    \item \textbf{Biểu đồ đoạn thẳng tần số tương đối ghép nhóm:} Nối các điểm có tọa độ $(c_i; f_i)$ bằng các đoạn thẳng liên tiếp để quan sát xu hướng phân phối hình chuông, lệch trái hoặc lệch phải của dữ liệu.
\end{itemize}

\paragraph{Mục 4: Tích hợp Năng lực số (Mã 5.1NC1b) \& Năng lực AI (Mã 11.C5.2) (03 phút)}

\begin{pedabox}[Năng lực số (Mã 5.1NC1b): Đếm tần số ghép nhóm bằng hàm FREQUENCY trên Excel]
    $\bullet$ \textbf{Thao tác tại phòng máy 35 máy tính:}\\
    1. Nhập dãy dữ liệu thô vào cột \texttt{A} (từ \texttt{A1} đến \texttt{A40}).\\
    2. Nhập các đầu mút phải của nhóm vào cột \texttt{B} (gọi là Bin range, ví dụ: $60, 80, 100, 120, 140$).\\
    3. Chọn dải ô ở cột tần số \texttt{C1:C5}, gõ công thức mảng: \texttt{=FREQUENCY(A1:A40, B1:B5)} rồi nhấn tổ hợp phím \texttt{Ctrl + Shift + Enter}. Excel sẽ tự động đếm chính xác số lượng giá trị cho từng nhóm trong tích tắc!\\
    4. Vào \texttt{Insert} $\to$ \texttt{Chart} $\to$ \texttt{Histogram} để hiển thị biểu đồ tần số tương đối tự động.
\end{pedabox}

\begin{aibox}[Tích hợp Năng lực AI (QĐ 2422/QĐ-BGDĐT --- Mã 11.C5.2): Thuật toán gom cụm (Clustering) trong AI]
    $\bullet$ \textbf{Bản chất toán học tương đồng:}\\
    Trong khoa học trí tuệ nhân tạo (AI / Machine Learning), thuật toán \textbf{Gom cụm K-Means (K-Means Clustering)} hoạt động dựa trên nguyên lí tương tự việc ghép nhóm số liệu: Từ một khối lượng dữ liệu khổng lồ không được gắn nhãn (Unsupervised Learning), AI tự động tìm các tâm cụm (tương tự giá trị đại diện $c_i$) và gom các điểm dữ liệu gần tâm nhất vào một cụm (tương tự nhóm $[a_i; a_{i+1})$).\\
    $\bullet$ \textbf{Ứng dụng thực tế của AI:} Hệ thống đề xuất của YouTube, TikTok gom nhóm người dùng có cùng sở thích thời lượng xem video; xe tự hành AI gom nhóm các vật thể có cùng khoảng cách đo từ cảm biến LiDAR.
\end{aibox}

\subsubsection*{3. Hoạt động 3: Luyện tập (12 phút)}
\begin{itemize}
    \item \textbf{Bài tập thực hành:} Khảo sát thời gian xem ti vi mỗi ngày của $30$ người cao tuổi tại một khu dân cư thu được bảng số liệu sau (đơn vị: phút):
    \begin{center}
    \begin{tabular}{|cccccccccc|}
    \hline
    $40$ & $55$ & $65$ & $70$ & $85$ & $90$ & $95$ & $100$ & $110$ & $125$ \\
    $45$ & $60$ & $68$ & $75$ & $85$ & $90$ & $98$ & $105$ & $115$ & $130$ \\
    $50$ & $60$ & $70$ & $80$ & $88$ & $92$ & $100$ & $110$ & $120$ & $135$ \\
    \hline
    \end{tabular}
    \end{center}
    Hãy ghép nhóm mẫu số liệu trên thành $5$ nhóm có độ dài bằng nhau gồm: $[40; 60)$, $[60; 80)$, $[80; 100)$, $[100; 120)$, $[120; 140]$.
    \begin{itemize}
        \item a) Lập bảng tần số và tần số tương đối ghép nhóm.
        \item b) Xác định giá trị đại diện của từng nhóm.
        \item c) Vẽ biểu đồ tần số tương đối ghép nhóm (Histogram).
    \end{itemize}
    \item \textit{Lời giải chi tiết:}\\
    a) \& b) Đếm số lượng quan sát theo quy tắc nửa khoảng $[a_i; a_{i+1})$:
    \begin{itemize}
        \item Nhóm $[40; 60)$: Gồm các giá trị $40 \le x < 60 \implies \{40, 45, 50, 55\} \implies m_1 = 4$.
        Giá trị đại diện: $c_1 = \dfrac{40 + 60}{2} = 50$.
        Tần số tương đối: $f_1 = \dfrac{4}{30} \cdot 100\% \approx 13{,}33\%$.
        \item Nhóm $[60; 80)$: Gồm các giá trị $60 \le x < 80 \implies \{60, 60, 65, 68, 70, 70, 75\} \implies m_2 = 7$.
        Giá trị đại diện: $c_2 = \dfrac{60 + 80}{2} = 70$.
        Tần số tương đối: $f_2 = \dfrac{7}{30} \cdot 100\% \approx 23{,}33\%$.
        \item Nhóm $[80; 100)$: Gồm các giá trị $80 \le x < 100 \implies \{80, 85, 85, 88, 90, 90, 92, 95, 98\} \implies m_3 = 9$.
        Giá trị đại diện: $c_3 = \dfrac{80 + 100}{2} = 90$.
        Tần số tương đối: $f_3 = \dfrac{9}{30} \cdot 100\% = 30{,}00\%$.
        \item Nhóm $[100; 120)$: Gồm các giá trị $100 \le x < 120 \implies \{100, 100, 105, 110, 110, 115\} \implies m_4 = 6$.
        Giá trị đại diện: $c_4 = \dfrac{100 + 120}{2} = 110$.
        Tần số tương đối: $f_4 = \dfrac{6}{30} \cdot 100\% = 20{,}00\%$.
        \item Nhóm $[120; 140]$: Gồm các giá trị $120 \le x \le 140 \implies \{120, 125, 130, 135\} \implies m_5 = 4$.
        Giá trị đại diện: $c_5 = \dfrac{120 + 140}{2} = 130$.
        Tần số tương đối: $f_5 = \dfrac{4}{30} \cdot 100\% \approx 13{,}34\%$.
    \end{itemize}

    \textbf{Bảng tần số và tần số tương đối ghép nhóm tổng hợp:}
    \begin{center}
    \begin{tabular}{|c|c|c|c|}
    \hline
    \textbf{Nhóm thời gian (phút)} & \textbf{Giá trị đại diện ($c_i$)} & \textbf{Tần số ($m_i$)} & \textbf{Tần số tương đối ($f_i$)} \\
    \hline
    $[40; 60)$ & $50$ & $4$ & $13{,}33\%$ \\
    $[60; 80)$ & $70$ & $7$ & $23{,}33\%$ \\
    $[80; 100)$ & $90$ & $9$ & $30{,}00\%$ \\
    $[100; 120)$ & $110$ & $6$ & $20{,}00\%$ \\
    $[120; 140]$ & $130$ & $4$ & $13{,}34\%$ \\
    \hline
    \textbf{Tổng cộng} & --- & $n = 30$ & $100{,}00\%$ \\
    \hline
    \end{tabular}
    \end{center}

    c) \textbf{Mô phỏng đồ họa biểu đồ Histogram tần số tương đối:}
    \begin{center}
    \begin{tikzpicture}[scale=0.85]
        \begin{axis}[
            ybar,
            bar width=1.1cm,
            ymin=0, ymax=35,
            ylabel={Tần số tương đối $f_i$ (\%)},
            xlabel={Thời gian (phút)},
            symbolic x coords={[40; 60), [60; 80), [80; 100), [100; 120), [120; 140]},
            xtick=data,
            nodes near coords,
            nodes near coords align={vertical},
            grid=both,
            grid style={dashed, gray!30},
            width=12cm,
            height=6.5cm
        ]
            \addplot[fill=teal!60, draw=teal!90] coordinates {
                ([40; 60), 13.3)
                ([60; 80), 23.3)
                ([80; 100), 30.0)
                ([100; 120), 20.0)
                ([120; 140], 13.4)
            };
        \end{axis}
    \end{tikzpicture}
    \end{center}
\end{itemize}

\subsubsection*{4. Hoạt động 4: Vận dụng thực tiễn (06 phút)}
\textbf{Chủ đề:} \textit{Phân tích phổ điểm thi thử tốt nghiệp THPT môn Toán}
\begin{itemize}
    \item \textbf{a) Mục tiêu:} Giúp học sinh hiểu cách ngành giáo dục và các trường đại học sử dụng mẫu số liệu ghép nhóm và biểu đồ phân phối điểm để đánh giá chất lượng dạy học và xác định điểm chuẩn trúng tuyển.
    \item \textbf{b) Nội dung:} Một trường THPT khảo sát điểm thi thử môn Toán của $500$ học sinh khối 12 được chia thành các nhóm điểm: $[0; 2)$, $[2; 4)$, $[4; 6)$, $[6; 8)$, $[8; 10]$. Biết rằng tần số tương đối của các nhóm lần lượt là: $2\%$, $10\%$, $45\%$, $33\%$, $10\%$.
    \begin{itemize}
        \item Có bao nhiêu học sinh đạt điểm khá - giỏi (từ $6{,}0$ điểm trở lên)?
        \item Hãy nhận xét về mức độ phân bố điểm số của học sinh nhà trường.
    \end{itemize}
    \item \textbf{c) Sản phẩm:}\\
    \textit{Lời giải:}\\
    - Tần số tương đối của học sinh đạt điểm từ $6{,}0$ điểm trở lên gồm nhóm $[6; 8)$ và $[8; 10]$:\\
    $f_{\ge 6} = 33\% + 10\% = 43\%$.\\
    Số học sinh đạt từ $6{,}0$ điểm trở lên là:
    $$N_{\ge 6} = 500 \cdot 43\% = 500 \cdot 0{,}43 = 215\text{ (học sinh)}$$
    - Nhận xét: Điểm số tập trung chủ yếu ở nhóm trung bình - khá $[4; 6)$ chiếm tới $45\%$. Tỉ lệ học sinh dưới trung bình ($< 5$ điểm) thấp, trong khi tỉ lệ khá giỏi đạt $43\%$, cho thấy đề thi có độ phân hóa tốt và chất lượng học tập của học sinh ở mức ổn định.
    \item \textbf{d) Tổ chức thực hiện:} GV nêu bài toán $\to$ HS tính toán nhanh $\to$ GV tổng kết ý nghĩa của việc trực quan hóa dữ liệu lớn trong đời sống.
\end{itemize}

\newpage

\section*{IV. PHỤ LỤC HỌC LIỆU BỔ TRỢ}

\subsection*{Phụ lục 1: Phiếu học tập số 1 --- Kĩ thuật xác định yếu tố nhóm số liệu}
\noindent\textbf{Họ và tên:} .................................................................... \textbf{Lớp:} 11A....... \textbf{Nhóm:} ..........

\begin{pedabox}[Nhiệm vụ 1: Hoàn thành thông tin định nghĩa nhóm số liệu]
    $\bullet$ Nhóm số liệu $[a; b)$ có:\\
    - Đầu mút trái là: ....................................... Đầu mút phải là: .......................................\\
    - Độ dài của nhóm: $h = $ .................................................................................................\\
    - Giá trị đại diện của nhóm: $c = $ ..................................................................................\\[0.3em]
    $\bullet$ Cho nhóm $[150; 160)$. Các số sau đây có thuộc nhóm hay không (Điền \textbf{Có} hoặc \textbf{Không}):\\
    - Số $150$: .................... \quad - Số $159{,}5$: .................... \quad - Số $160$: ....................
\end{pedabox}

\subsection*{Phụ lục 2: Phiếu học tập số 2 --- Thực hành phân nhóm dữ liệu thực tế}
\noindent\textbf{Họ và tên:} .................................................................... \textbf{Lớp:} 11A....... \textbf{Nhóm:} ..........

\begin{pedabox}[Nhiệm vụ 2: Điền bảng tần số và tần số tương đối ghép nhóm]
Khảo sát cân nặng (kg) của $20$ học sinh: $45, 48, 50, 52, 55, 47, 49, 53, 56, 58, 60, 62, 51, 54, 57, 59, 46, 50, 53, 55$.\\
Hãy hoàn thành bảng sau với $4$ nhóm: $[45; 50), [50; 55), [55; 60), [60; 65]$:
\begin{center}
\begin{tabularx}{0.9\linewidth}{|c|X|X|X|}
\hline
\textbf{Nhóm cân nặng} & \textbf{Giá trị đại diện} & \textbf{Tần số ($m_i$)} & \textbf{Tần số tương đối (\%)} \\
\hline
$[45; 50)$ & & & \\
$[50; 55)$ & & & \\
$[55; 60)$ & & & \\
$[60; 65]$ & & & \\
\hline
\textbf{Tổng} & --- & \textbf{20} & \textbf{100\%} \\
\hline
\end{tabularx}
\end{center}
\end{pedabox}

\subsection*{Phụ lục 3: Bảng Rubric đánh giá năng lực tích hợp trong Bài 8}

\begin{center}
\begin{tabularx}{\linewidth}{|p{3.2cm}|X|X|X|X|}
\hline
\textbf{Tiêu chí} & \textbf{Mức 1 (Chưa đạt)} & \textbf{Mức 2 (Đạt)} & \textbf{Mức 3 (Khá)} & \textbf{Mức 4 (Tốt)} \\
\hline
\textbf{1. Kĩ năng ghép nhóm số liệu} & Nhầm lẫn quy tắc lấy dấu bằng ở đầu mút; đếm sót hoặc đếm trùng lặp. & Xác định đúng các phần tử thuộc nửa khoảng $[a; b)$; tính đúng giá trị đại diện. & Lập chuẩn xác bảng tần số và bảng tần số tương đối ghép nhóm; tính tổng đúng $100\%$. & Vẽ thành thạo biểu đồ Histogram cột liền kề chuẩn tỉ lệ; đọc hiểu sâu sắc ý nghĩa phân bố. \\
\hline
\textbf{2. Ứng dụng số (Mã 5.1NC1b)} & Chưa biết cách mở bảng tính hoặc nhập liệu sai cột. & Nhập đúng số liệu thô vào cột Excel trên phòng máy. & Áp dụng thành thạo hàm mảng \texttt{FREQUENCY} để tự động đếm tần số nhóm. & Tự thiết lập biểu đồ Histogram tự động và chỉnh sửa định dạng đồ họa số liệu chuyên nghiệp. \\
\hline
\textbf{3. Năng lực AI (Mã 11.C5.2)} & Chưa hiểu khái niệm gom cụm trong máy học. & Nêu được tên thuật toán gom cụm (K-Means) trong trí tuệ nhân tạo. & Phân tích được sự tương đồng giữa chia nhóm số liệu thống kê và gom cụm dữ liệu AI. & Đánh giá sâu sắc vai trò của việc tiền xử lý và gom nhóm dữ liệu lớn trong huấn luyện mô hình AI. \\
\hline
\end{tabularx}
\end{center}

\end{document}
